\documentclass[11pt]{article}
\usepackage[a4paper,margin=21mm]{geometry}
\usepackage{fontspec}
\setmainfont{Latin Modern Roman}
\newfontfamily\CJKfont{Noto Serif CJK SC}
\XeTeXlinebreaklocale "zh"
\XeTeXlinebreakskip=0pt plus 1pt
\usepackage{amsmath,amssymb,amsthm,amscd,graphicx,color}
\usepackage{xurl}
\usepackage[unicode,hidelinks]{hyperref}
\allowdisplaybreaks
\setlength\emergencystretch{3em}
\numberwithin{equation}{section}

\numberwithin{equation}{section}

\usepackage{microtype}

\usepackage[all]{xy}

\usepackage{mathtools}
\usepackage{stackrel}

\usepackage{csquotes}

\newtheorem{thm}{Theorem}[section]
\newtheorem{lem}[thm]{Lemma}
\newtheorem{prop}[thm]{Proposition}
\newtheorem{cor}[thm]{Corollary}
\newtheorem{question}[thm]{Question}
\newtheorem{problem}[thm]{Problem}
\newtheorem{claim}[thm]{Claim}
\newtheorem{obs}[thm]{Observation}
\newtheorem{conj}[thm]{Conjecture}
\newtheorem{ass}[thm]{Assumption}

\theoremstyle{definition}
\newtheorem{defn}[thm]{Definition}
\newtheorem{Example}[thm]{Example}
\newtheorem{examples}[thm]{Examples}
\newtheorem{rem}[thm]{Remark}
\newtheorem{nota}[thm]{Notation}
\newtheorem{stancon}[thm]{Standard construction}

\newcommand{\N}{\mathbb{N}}
\newcommand{\Z}{\mathbb{Z}}
\newcommand{\Q}{\mathbb{Q}}
\newcommand{\R}{\mathbb{R}}
\newcommand{\C}{\mathbb{C}}

\newcommand{\cZ}{\mathcal{Z}}
\newcommand{\cO}{\mathcal{O}}
\newcommand{\cP}{\mathcal{P}}
\newcommand{\cF}{\mathcal{F}}
\newcommand{\cX}{\mathcal{X}}
\newcommand{\cY}{\mathcal{Y}}
\newcommand{\cA}{\mathcal{A}}
\newcommand{\cB}{\mathcal{B}}
\newcommand{\cU}{\mathcal{U}}
\newcommand{\cV}{\mathcal{V}}
\newcommand{\cW}{\mathcal{W}}
\newcommand{\cK}{\mathcal{K}}
\newcommand{\cL}{\mathcal{L}}
\newcommand{\cC}{\mathcal{C}}

\newcommand{\Ct}{\mathrm{C}}
\newcommand{\Cz}{\Ct_0}
\newcommand{\Cb}{\Ct_b}
\newcommand{\Ch}{\Ct_h}
\newcommand{\CzRed}{\Ct_0^{\mathrm{red}}}
\newcommand{\CRed}{\Ct^{\mathrm{red}}}

\newcommand{\fr}{\mathfrak{r}}
\newcommand{\fV}{\mathfrak{V}}
\newcommand{\fP}{\mathfrak{P}}

\newcommand{\LocComp}{\mathbf{LCH}}
\newcommand{\sigLocComp}{\boldsymbol\sigma\mathbf{LCH}}
\newcommand{\Comp}{\mathbf{CH^2}}
\newcommand{\sigComp}{\boldsymbol\sigma\mathbf{CH^2}}

\newcommand{\categoryfont}[1]{\mathsf{#1}}
\newcommand{\sLCH}[1][\Gamma]{\boldsymbol\sigma\categoryfont{LCH}_{#1}}
\newcommand{\sLCHz}[1][\Gamma]{\boldsymbol\sigma\categoryfont{LCH}_{#1}^2}
\newcommand{\sLCHp}[1][\Gamma]{\boldsymbol\sigma\categoryfont{LCH}_{#1}^+}
\newcommand{\sLCHzp}[1][\Gamma]{\boldsymbol\sigma\categoryfont{LCH}_{#1}^{+,2}}
\newcommand{\sCH}[1][\Gamma]{\boldsymbol\sigma\categoryfont{CH}_{#1}}
\newcommand{\sCHz}[1][\Gamma]{\boldsymbol\sigma\categoryfont{CH}_{#1}^2}
\newcommand{\LCH}[1][\Gamma]{\categoryfont{LCH}_{#1}}
\newcommand{\LCHz}[1][\Gamma]{\categoryfont{LCH}_{#1}^2}
\newcommand{\LCHp}[1][\Gamma]{\categoryfont{LCH}_{#1}^+}
\newcommand{\LCHzp}[1][\Gamma]{\categoryfont{LCH}_{#1}^{+,2}}
\newcommand{\CH}[1][\Gamma]{\categoryfont{CH}_{#1}}
\newcommand{\CHz}[1][\Gamma]{\categoryfont{CH}_{#1}^2}

\newcommand{\sCs}[1][\Gamma]{\boldsymbol\sigma\categoryfont{C}^*_{#1}}
\newcommand{\csCs}[1][\Gamma]{\categoryfont{c}\boldsymbol\sigma\categoryfont{C}^*_{#1}}

\newcommand{\CGCBG}[1][\Gamma]{\categoryfont{CGBG}_{#1}}
\newcommand{\CGCBGz}[1][\Gamma]{\categoryfont{CGBG}^2_{#1}}

\newcommand{\coarseStr}{\mathcal{E}}
\newcommand{\topology}{\mathcal{T}}
\newcommand{\Diag}[1][]{{\Delta_{#1}}}
\newcommand{\swapcross}{\mathbin{\bar\times}}
\newcommand{\HilbertMod}{\mathfrak{H}}
\newcommand{\Paschke}{\mathfrak{D}}
\newcommand{\PaschkeRed}{\mathfrak{D}^{\textrm{red}}}

\DeclareMathOperator{\im}{im}
\DeclareMathOperator{\id}{id}
\DeclareMathOperator{\diam}{diam}
\DeclareMathOperator{\dist}{dist}
\DeclareMathOperator{\Var}{Var}
\DeclareMathOperator{\codim}{codim}
\DeclareMathOperator{\supp}{supp}
\DeclareMathOperator{\Pen}{Pen}
\DeclareMathOperator{\Cone}{Cone}
\DeclareMathOperator{\Ad}{Ad}
\DeclareMathOperator{\Prop}{prop}

\newcommand{\ev}[1]{\mathop{\textrm{ev}_{#1}}}
\newcommand{\elltwo}{\ell^2}

\newcommand{\maxtensor}{\mathbin{{\otimes}_{\mathrm{max}}}}
\newcommand{\maxcrossed}{\mathbin{{\rtimes}_{\mathrm{max}}}}
\newcommand{\redcrossed}{\mathbin{{\rtimes}_{\mathrm{red}}}}


\newcommand{\Homol}{\mathrm{E}}
\newcommand{\Cohom}{\mathrm{E}}
\newcommand{\HomolX}{\mathrm{EX}}
\newcommand{\CohomX}{\mathrm{EX}}

\newcommand{\bound}[1][0]{\partial^{{\textsc{\romannumeral #1}}}}
\newcommand{\cobound}[1][0]{\delta_{{\textsc{\romannumeral #1}}}}

\newcommand{\HX}{\mathrm{HX}}
\newcommand{\CX}{\mathrm{CX}}
\newcommand{\KX}{\mathrm{KX}}
\newcommand{\KXRed}{\widetilde{\mathrm{K}}\mathrm{X}}
\newcommand{\K}{\mathrm{K}}
\newcommand{\Ktop}{\mathrm{K}^{\mathrm{top}}}
\newcommand{\KK}{\mathrm{KK}}
\newcommand{\HA}{\mathrm{HA}}
\newcommand{\CA}{\mathrm{CA}}
\newcommand{\CZ}{\mathrm{CZ}}
\newcommand{\CK}{\mathrm{CK}}
\newcommand{\EE}{\mathrm{E}}
\newcommand{\HS}{\mathrm{HS}}
\newcommand{\CS}{\mathrm{CS}}
\newcommand{\lf}{\mathrm{lf}}
\newcommand{\cs}{\mathrm{c}}


\newcommand{\textCstar}{\ensuremath{\mathrm{C}^*\!}}


\newcommand{\Roe}{\mathrm{C}^*}
\newcommand{\PseudoLoc}{\mathrm{D}^*}
\newcommand{\Loc}{\mathrm{C}^*_{\mathrm{L},\Gamma}}
\newcommand{\PseudoLocLoc}{\mathrm{D}^*_{\mathrm{L},\Gamma}}
\newcommand{\sHigCom}{\overline{\mathfrak{c}}}
\newcommand{\sHigCor}{\mathfrak{c}}
\newcommand{\sHigComRed}{\overline{\mathfrak{c}}^{\mathrm{red}}}
\newcommand{\sHigCorRed}{\mathfrak{c}^{\mathrm{red}}}
\newcommand{\sHigFuAlg}{\overline{\mathfrak{b}}}
\newcommand{\sHigFuAlgRed}{\overline{\mathfrak{b}}^{\mathrm{red}}}
\newcommand{\VanInfFuAlg}{\mathrm{B}_0}
\newcommand{\HigFuAlg}{\mathrm{B}_h}
\newcommand{\Kom}{{\mathfrak{K}}}
\newcommand{\Lin}{{\mathfrak{B}}}
\newcommand{\multiplier}{{\mathcal{M}}}
\newcommand{\doubleplier}{{\mathcal{D}}}
\newcommand{\Asymp}{\mathfrak{A}}
\newcommand{\admissible}{{\mathfrak{C}}}

\newcommand{\sfP}{\mathsf{P}}
\newcommand{\sfD}{\mathsf{D}}

\newcommand{\blank}{-}
\newcommand{\coarsecomp}{\Pi_0}
\newcommand{\myvdots}{{\vphantom{(}\raisebox{-0.3ex}[0pt][0pt]{\vdots}}}



\makeatletter\newlabel{sec:Sigmastuff}{{5}{}{Original source locator}{}{}}\makeatother
\begin{document}
\begin{center}\Large The canonical localization-algebra map is an equivariant E-theory K-homology isomorphism for proper cocompact isometric actions with separable coefficients\end{center}
\noindent\textbf{Stable ID:} 1250\quad\textbf{Author:} Christopher Wulff\par
\noindent\textbf{Work:} Equivariant Coarse (Co-)Homology Theories\par
\noindent\textbf{Source:} \url{https://sigma-journal.com/2022/057/}\par
\noindent\textbf{Version:} Official SIGMA 18 (2022) article 057, DOI 10.3842/SIGMA.2022.057; journal PDF and native TeX acquired 2026-09-30, exact bytes pinned by SHA256\par
\noindent\textbf{Rights:} CC BY 4.0. Authors retain copyright. \url{https://creativecommons.org/licenses/by/4.0/}. Reuse and typesetting adaptation; original language and mathematical content preserved.\par
\noindent\textbf{Difficulty (editorial):} {\CJKfont H2: The proof compares concrete finite-propagation localization operators with equivariant KK cycles and then identifies the resulting transformation with the canonical map to E-theory. It constructs the projection-valued representative of a localized unitary class and checks the coefficient-natural transformation on the identity cycle using an explicit zero-propagation projection. This verifies the particular asymptotic-morphism map, rather than merely asserting that the two groups are abstractly isomorphic.}\par
\medskip\noindent\textbf{Editorial boundary note (not author text):} Selected absolute Theorem6.8 canonical Delta isomorphism for proper metric X with countable discrete Gamma acting properly, isometrically and cocompactly, and separable Gamma-C*-coefficients D. Full printed broader theorem/author historical footnote remain context; noncocompact truncation, nonseparable extension and relative consequence receive no count. Complete Hilbert-module/Roe/local-compactness/propagation definitions, actual localization algebra, delta asymptotic morphism and exact Delta composition are retained. The previously established cocompact Paschke/localization isomorphisms stay author-cited prerequisites. Full current even-degree unitary/projection representative and coefficient-natural identity-cycle computation, followed by suspension to odd degree, are unchanged. Only the separate unselected noncocompact modification paragraph2361–2364 is excluded between byte-bound proof blocks; no proof repair or generated substitution. Original footnote14–17 numbering is restored by wrapper-only counter reset.\par
\noindent\textbf{External original reference sec:Sigmastuff:} Original sigma-space general theory section5, native1673–2095, contextual reference only; selected absolute proper metric case uses defined E-theory directly\par
\medskip\hrule\medskip\noindent\textbf{Author text begins below.}\par
\setcounter{section}{4}\setcounter{subsection}{3}
% Original source lines 1442--1476
\subsection[K-theory of the Roe algebra]{$\boldsymbol{\K}$-theory of the Roe algebra}\label{sec:RoeAlg}

The Roe algebra is a \textCstar-algebra $\Roe(X)$ that can be constructed for every coarse space $X$ (cf.\ \cite[Section~4.4]{RoeCoarseGeometry}), but it is usually only dealt with in the case of proper metric spaces. As most of the references only consider the latter special case, we shall stick to it.

Let us briefly recall the broad outlines of its construction as presented, e.g., in \cite[Section~6.3]{HigRoe}. Starting with a representation $\rho$ of $\Cz(X)$ on a separable Hilbert space $H$ which is non-degenerate (meaning that $\rho(\Cz(X))H$ is dense in $H$), the Roe algebra $\Roe(X,\rho)$ is defined as the norm closure of the $*$-subalgebra of $\Lin(H)$ consisting of all operators $T$ satisfying the following two properties:\looseness=1
\begin{enumerate}\itemsep=0pt
\item[$\bullet$] $T$ is called \emph{locally compact}, if $\rho(f)T,T\rho(f)$ are compact operators on $H$ for all $f\in\Cz(X)$.
\item[$\bullet$] $T$ has \emph{finite propagation}, if there is $R\geq 0$ such that $\rho(f)T\rho(g)=0$ whenever the distance between the supports of $f,g\in\Cz(X)$ is at least $R$. The infimum of all $R\geq 0$ with this property is called the \emph{propagation} $\Prop(T)$ of $T$.
\end{enumerate}
This \textCstar-algebra depends on the choice of $H$ and $\rho$, but if the representation is big enough in a~cer\-tain sense (it is \enquote{ample}, i.e., it is non-degenerate and $\rho^{-1}(\Kom(H))=\{0\}$), then its $\K$-theory $\K_*(\Roe(X,\rho))$ is independent from the choices up to canonical isomorphism.
Therefore, the representation $\rho$ is usually omitted from the notation with the understanding that one suitable representation has been chosen for each $X$.

Now, there are two well-known generalizations of the Roe algebra. The first is the equivariant Roe algebra in the case that a countable discrete group $\Gamma$ acts properly and isometrically on the coarsely proper metric space $X$. Here one chooses in addition to the non-degenerate representation~$\rho$ also a representation $u$ of the group $\Gamma$ on $H$ by unitaries such that $u(\gamma)\rho(f)u(\gamma)^*=\rho(\gamma\cdot f)$ for all $\gamma\in\Gamma$ and $f\in\Cz(X)$. The triple $(H,\rho,u)$ is called a $\Gamma$-$X$-module. Then the equivariant Roe algebra $\Roe_\Gamma(X,\rho,u)$ is the norm closure of the \textCstar-algebra of all locally compact operators of finite propagation in $\Lin(H)$ that are equivariant with respect to $u$. Again, the $\K$-theory groups $\K_*(\Roe_\Gamma(X,\rho,u))$ are independent up to canonical isomorphism from the $\Gamma$-$X$-module $(H,\rho,u)$ if the latter has been chosen sufficiently large and then abbreviated by $\K_*(\Roe_\Gamma(X))$.
A detailed exposition about equivariant Roe algebras can be found in \cite{WillettYuHigherIndexTheory}.

The second generalization is to introduce a (possibly graded) coefficient \textCstar-algebra $D$ into the Roe algebra, cf.\ \cite{HankePapeSchick,HigsonPedersenRoe,WulffTwisted}. This is done by replacing the Hilbert space $H$ by a right Hilbert module $\HilbertMod$ over $D$ and the non-degenerate representation $\rho$ is by adjointable operators on $\HilbertMod$.
The $\K$-theory of the resulting \textCstar-algebra $\Roe(X,\rho,\HilbertMod;D)$ is again unique up to canonical isomorphism if $\rho$ and $\HilbertMod$ have been chosen sufficiently large and in this case $\rho$ will be ommited from the notation. Here, sufficient largeness means that $\HilbertMod\cong H\otimes D$ and $\rho=\rho'\otimes \id_D$ for an ample representation~$\rho'$ on a separable Hilbert space $H$ \cite[Definition 4.5]{HigsonPedersenRoe}.


It is possible to unify the two generalizations to equivariant Roe algebras $\Roe_\Gamma(X,\rho,u,\HilbertMod;D)$ with coefficients in a (graded) $\Gamma$-\textCstar-algebra $D$.
In this case, $u$ also has to be a representation by adjointable operators on $\HilbertMod$ which is compatible with the $\Gamma$-action on $D$ in the sense that $u(\gamma)(\xi d)=(u(\gamma)\xi) (\gamma d)$ for all $\xi\in\HilbertMod$, $d\in D$ and $\gamma\in\Gamma$.
Again, we obtain unambiguous $\K$-theory groups $\K_*(\Roe_\Gamma(X;D))$ if we start with sufficiently large $H,\rho,u$ as in the equivariant case and then simply use their tensor product with $D$, i.e., $\HilbertMod\coloneqq H\otimes D$ with obvious induced representations of $\Cz(X)$ and $\Gamma$.




The proofs of the following statements can be found in the above-mentioned sources in the special cases just described, but the proofs generalize to the general case of equivariant Roe-algebras with coefficients.

If $A\subset X$ is a $\Gamma$-invariant closed subspace of the proper metric $\Gamma$-space $X$, then $\Roe_\Gamma(X;D)$ contains an ideal $\Roe_\Gamma(A\subset X;D)$ whose $\K$-theory is canonically isomorphic to the $\K$-theory of $\Roe_\Gamma(A;D)$.
It is the norm closure of all $T\in\Roe_\Gamma(X;D)$ which are supported in an $R$-neighborhood of $A$ for some $R\geq 0$, that is, all those $T$ for which $\rho(f)T=T\rho(f)=0$ for all $f\in\Cz(X)$ with $\dist(\supp(f),A)\geq R$.
The quotient \textCstar-algebra $\Roe_\Gamma(X,A;D)\coloneqq \Roe_\Gamma(X;D)/\Roe_\Gamma(A\subset X;D)$ is called the \emph{$\Gamma$-equivariant relative Roe algebra}. We have $\Roe_\Gamma(\varnothing\subset X;D)=0$ and hence $\Roe_\Gamma(X,\varnothing;D)=\Roe_\Gamma(X;D)$. The long exact sequences
\[
\dots\to\K_p(\Roe_\Gamma(A;D))\to \K_p(\Roe_\Gamma(X;D))\to \K_p(\Roe_\Gamma(X,A;D))\to \K_{p-1}(\Roe_\Gamma(A;D))\to \cdots
\]
are immediate.

\setcounter{section}{6}\setcounter{subsection}{3}\setcounter{thm}{6}\setcounter{equation}{1}\setcounter{footnote}{13}
% Original source lines 2282--2359
\subsection[Coarse K-homology and assembly]{Coarse $\boldsymbol{\K}$-homology and assembly}\label{sec:Assembly}

Using $\Gamma$-equivariant $\EE$-theory,\footnote{Of course, $\KK$-theory could also be used. In that case, a different description of the assembly maps would need to be used and hence completely new proofs of our results would be required. However, $\KK$-theory suffers from similar technical shortcomings as $\EE$-theory and therefore no advantage over our approach is expected. In the end, $\EE$-theory was primarily selected because of the author's greater familiarity with this theory.} we define the $\Gamma$-equivariant $\K$-homology with coefficients in a~separable $\Gamma$-\textCstar-algebra $D$ as the functors
\[
\K^\Gamma_*(\blank;D)\coloneqq \EE^\Gamma_*(\Cz(\blank);D).
\]
Note that this definition requires the function algebras $\Cz(\blank)$ to be separable, so we only consider these functors on the full subcategory of $\LCHp$ of all \emph{second countable} locally compact $\Gamma$-spaces.
In practice, this is only a minor restriction to the theory of Section~\ref{sec:Sigmastuff}, because we are mainly interested in two special cases: First, Rips complexes of discretizations of countably generated proper isocoarse $\Gamma$-spaces of bornologically bounded geometry are always second countable and, second, their compactification with a Higson dominated corona is second countable, if the corona is second countable, i.e., a compact metrizable space. The latter is an ubiquitous assumption in applications.


Due to the well-known properties of bivariant $\K$-theory (cf.\ \cite[Chapter~6]{GueHigTro}), the functors $\K^\Gamma_*(\blank;D)$ are $\Gamma$-equivariant single-space homology theories on the above-mentioned category. Furthermore, as direct limits preserve exact sequences, they can be extended to $\Gamma$-equivariant single-space homology theories on the full subcategory of $\sLCHp$ consisting of all second countable $\sigma$-locally compact $\Gamma$-spaces $\cX=\bigcup_{n\in\N}X_n$ by
\(\K^\Gamma_*(\cX;D)\coloneqq\lim_{n\to\infty}\K^\Gamma_*(X_n;D)\).


We want to compare the resulting transgression maps with the coarse assembly maps $\mu$: $\KX^\Gamma_*(X;D)\to\K_*(\Roe_\Gamma(X;D))$. The definition of the latter is based upon other pictures of $\K$-homology.
We choose the one using Yu's localization algebras \cite{YuLocalization} (see also \cite{QiaoRoe} and \cite{EngelWulffZeidler}) for our discussion.
A lot about localization algebras has been proven in the comprehensive book~\cite{WillettYuHigherIndexTheory}, but note that the ``localized Roe algebras'' considered there are slightly bigger than Yu's original localization algebras and hence a bit of care has to be taken when applying their results.

Following up our Section~\ref{sec:RoeAlg} about Roe algebras, let $X$ be a proper metric space equipped with a proper isometric $\Gamma$-action. We assume that sufficiently large $\Gamma$-$X$-module $(H,\rho,u)$ has been chosen and let $\HilbertMod\coloneqq H\otimes D$ be equipped with the induced representations. Recall that we defined $\Roe_\Gamma(X;D)$ as the norm closure in $\Lin(\HilbertMod)$ of the $*$-subalgebra of all $\Gamma$-equivariant locally compact operators of finite propagation, and if $A$ is a subspace of $X$, then $\Roe_\Gamma(A\subset X;D)$ denoted the ideal generated by all operators which are supported in an $R$-neighborhood of $A$ for some $R>0$ and $\Roe_\Gamma(X,A;D)\coloneqq \Roe_\Gamma(X;D)/\Roe_\Gamma(A\subset X;D)$.


\begin{defn}
%\label{defn:localizationalgebra}
The \emph{localization algebra} $\Loc(X;D)$ is the norm closure of the sub-$*$-algebra of $\Cb([1,\infty),\Roe_\Gamma(X;D))$ of all bounded uniformly continuous functions $L\colon[1,\infty)\to\Roe_\Gamma(X;D)$ such that the propagation of $L(t)$ is finite for all $t\geq 1$ and converges to zero as $t\to\infty$. Furthermore, we denote by $\Loc(A\subset X;D)\subset \Loc(X;D)$ the ideal generated by all $L$ for which there is a~function $\varepsilon\colon[1,\infty)\to(0,\infty)$ with $\varepsilon(t)\xrightarrow{t\to\infty}0$ such that $L(t)$ is supported in an $\varepsilon(t)$-neighborhood of $A$ for all $t\geq 1$, and we define $\Loc(X,A;D)\coloneqq \Loc(X;D)/\Loc(A\subset X;D)$.
\end{defn}

An obvious variant of \cite[Lemma~6.3.6]{WillettYuHigherIndexTheory}, which can be proven with the same methods, says that the inclusion $A\subset X$ induces via the theory of covering isometries an isomorphism $\K_*(\Loc(A;D))\cong\K_*(\Loc(A\subset X;D))$. Therefore we obtain a long exact sequence
\begin{align}
\cdots&\to\K_*(\Loc(A;D))\to \K_*(\Loc(X;D))\to \K_*(\Loc(X,A;D))\nonumber
\\&\to \K_{*-1}(\Loc(A;D))\to\cdots.\label{eq:localizationLES}
\end{align}


The relation between localization algebras and our definition of $\K$-homology is the following. As $\Kom(\HilbertMod)\cong D\otimes\Kom(H)$, it is straightforward to verify (cf.\ \cite[Corollary~4.2]{QiaoRoe}) that
\begin{align*}
\delta\colon\ \Loc(X;D)\maxtensor\Cz(X) &\to\frac{\Cb([1,\infty),D\otimes \Kom(H))}{\Cz([1,\infty),D\otimes\Kom(H))},
\\
L\otimes f&\mapsto [t\mapsto L(t)\circ (\rho(f)\otimes\id_D)]
\end{align*}
defines an $\Gamma$-equivariant asymptotic morphism. It vanishes on the ideal $\Loc(A\subset X;D)\maxtensor\Cz(X\setminus A)$ and hence we also obtain an induced $\Gamma$-equivariant asymptotic morphism
\[
\delta\colon\ \Loc(X,A;D)\maxtensor\Cz(X\setminus A) \to\frac{\Cb([1,\infty),D\otimes \Kom(H))}{\Cz([1,\infty),D\otimes\Kom(H))}.
\]
Both of them define canonical $\EE^\Gamma$-theory elements in $\EE^\Gamma_0(\Loc(X;D)\maxtensor\Cz(X),D)$ and \linebreak $\EE^\Gamma_0(\Loc(X,A;D)\maxtensor\Cz(X\setminus A),D)$, respectively, which we shall also denote by the letter~$\delta$.\footnote{Here we need $\EE^\Gamma$-theory also for non-separable \textCstar-algebras. It can be defined in an ad hoc way as in \cite[Section~1]{WulffTwisted}, or even simpler by letting $\EE^\Gamma_*(A,B)$ for possibly non-separable $A$, $B$ be obtained from the groups $\EE^\Gamma_*(A',B')$ by taking the direct limit over all separable $B'\subset B$ and the inverse limit over all separable $A'\subset A$. The relevant maps considered in the remainder of this section will be the same for both versions.
Note that these generalizations retain the composition and tensor products of $\EE^\Gamma_*$-theory, but be aware that they are simply too naive to preserve the long exact sequences. }
We obtain the composed group homomorphism
\begin{align*}
\Delta\colon\ \K_*(\Loc(X;D))&\cong\EE_*(\C,\Loc(X;D))
\\&\xrightarrow{\blank\otimes\id_{\Cz(X)}}\EE^\Gamma_*(\Cz(X),\Loc(X;D)\maxtensor\Cz(X))
\\&\xrightarrow{\delta\circ\blank}\EE^\Gamma_*(\Cz(X),D)=\K^\Gamma_*(X;D)
\end{align*}
and similarily $\Delta\colon \K_*(\Loc(X,A;D))\to\K^\Gamma_*(X,A;D)$ in the relative case.
By exploiting the homological properties of domains and targets of the maps $\Delta$ (in particular that the domains are functorial under uniformly continuous coarse $\Gamma$-maps via the theory of covering isometries, that the maps $\Delta$ are natural under such maps and that the connecting homomorphisms are in both cases constructed via mapping cones) we see that the maps $\Delta$ map \eqref{eq:localizationLES} to the long exact sequence of $\K^\Gamma_*(\blank,D)$.

\begin{thm}\label{thm:Khomologydefinitions}
The maps $\Delta$ are isomorphisms in the absolute case for every proper metric space~$X$ equipped with a proper isometric $\Gamma$-action and hence by a five lemma argument also in the relative case.\footnote{Instead of Theorem~\ref{thm:Khomologydefinitions}, preprint versions of this article contained the ``Assumption 6.8'' that $\Gamma$ and $D$ have been chosen for which the statement of this theorem holds. Thanks to a referee's hint, we can now prove it in full generality.}
\end{thm}

Proofs of some special cases of the theorem have been known for a long time.
In the non-equivariant case ($\Gamma=1$) without coefficients ($D=\C$), one proof is sketched in \cite[Construction~6.7.5, Exercises~B.3.2 and~B.3.3]{WillettYuHigherIndexTheory}, albeit for a different version of the localization algebra. Note, however, that this proof relies solely on homological properties and a comparison on one-point spaces. Therefore, it goes through for original version of the localization algebras, too, and even with arbitrary coefficient \textCstar-algebras $D$.

Another proof for the case $\Gamma=1$ and $D=\C$ can be found in \cite[Proposition~4.3]{QiaoRoe} and relies on Paschke duality. The author is indebted to a referee for pointing out a version of Paschke duality for the equivariant $\KK$-theory groups $\KK^\Gamma_*(\Cz(X),D)$ which allows to prove the theorem in full generality.
Nevertheless, as localization algebras and $\EE$-theory are defined in a very similiar manner, it would be interesting to know if the theorem can be proven directly without taking the detour via Paschke duality.

\begin{proof}
It suffices to prove the theorem for separable $D$, because the general case follows by taking direct limits over all separable sub-\textCstar-algebras.
In this case, we can make use of the Paschke duality for equivariant $\KK$-theory that was discussed in \cite[Appendix~B]{GueWillYu_DyncomplKtheory}.

An adjointable operator $T$ on $\HilbertMod$ is called pseudolocal if the commutators $\rho(f)T-T\rho(f)$ are compact for all $f\in \Cz(X)$.
Let $\PseudoLoc_\Gamma(X;D)$ denote the sub-\textCstar-algebra of $\Lin(\HilbertMod)$ generated by the $\Gamma$-equivariant pseudolocal operators that have finite propagation and let $\PseudoLocLoc(X;D)$ denote the sub-\textCstar-algebra of $\Cb([1,\infty),\PseudoLoc_\Gamma(X;D))$ of all bounded uniformly continuous functions $L\colon[1,\infty)\to\PseudoLoc(X;D)$ such that the propagation $L(t)$ is finite for all $t\geq 1$ and converges to zero as $t\to\infty$.\looseness=-1

We are free to choose a very specific sufficiently large $\HilbertMod=H\otimes D$, namely $\HilbertMod=\HilbertMod_D\coloneqq H'\otimes \ell^2(\N)\otimes\ell^2(\Gamma)\otimes D$ for a non-degenerate $\Gamma$-$X$-module $(H',\rho',u')$.
It will be equipped with $\Gamma$-action $\varepsilon$ which acts diagonally on the first, third and fourth factor and the representation $\pi$ of $\Cz(X)$ which acts on the first factor.
For this choice, \cite[Appendix B]{GueWillYu_DyncomplKtheory} contains the proof of isomorphisms
\begin{gather*}
\KK_p^\Gamma(\Cz(X),D)\cong\K_{p+1}\biggl(\frac{\PseudoLoc_\Gamma(X;D)}{\Roe_\Gamma(X;D)}\biggr) \cong\K_{p+1}\biggl(\frac{\PseudoLocLoc(X;D)}{\Loc(X;D)}\biggr)\cong\K_p(\Loc(X;D))
\end{gather*}
for separable $\Gamma$-\textCstar-algebras $D$ and proper metric spaces $X$ equipped with a proper cocompact isometric $\Gamma$-action.


% Original source lines 2366--2382
Let us briefly recall the definition of the abovementioned isomorphisms in degree $p=0$.
According to the analogue of \cite[Lemma B.7]{GueWillYu_DyncomplKtheory} in even degree\footnote{The cited lemma only considers the odd case, but the even case can be proven completely analogously.}, every class in $\KK_*^\Gamma(\Cz(X),D)$ can be represented by a cycle of the form $\left(\HilbertMod_D\oplus\HilbertMod_D,\left(\begin{smallmatrix}0&U^*\\U&0\end{smallmatrix}\right)\!,\epsilon\oplus\epsilon,\pi\oplus\pi\right)$ with $U\in \PseudoLoc_\Gamma(X;D)$.
Then $U$ is a unitary modulo $\Roe_\Gamma(X;D)$ and the Paschke duality isomorphism is defined by mapping this cycle to the class represented by $U$. The second isomorphism then takes it to the class represented by a function $t\mapsto U_t$ which represents a unitary in $\PseudoLocLoc(X;D)/\Loc(X;D)$ and such that $U_0$ agrees with $U$ modulo $\Roe_\Gamma(X;D)$. Such a function can be constructed from the constant function $U$ by applying the procedure from the proof of \cite[Lemma~B.15]{GueWillYu_DyncomplKtheory}.
The third isomorphism then maps it to the class represented by the difference of the projection valued functions $t\mapsto P_t$ and $t\mapsto Q$ defined by
\[
P_t\coloneqq \begin{pmatrix}U_tU_t^*&U_t(\id-U_t^*U_t)^{\frac12} \\U_t^*(\id-U_tU_t^*)^{\frac12}&\id-U_t^*U_t\end{pmatrix}\qquad\text{and}\qquad Q=\begin{pmatrix}\id&0\\0&0\end{pmatrix}\!,
\]
cf.\ \cite[Proposition~4.8.10]{HigRoe}.

In order to prove the claim in even degrees, it now suffices to show that $\Delta$ composed with these isomorphisms is the canonical natural isomorphism $\KK_0^\Gamma(\Cz(X),D)\xrightarrow{\cong}\EE_0^\Gamma(\Cz(X),D)$.
The latter is defined via \cite[Theorem~2.2]{Thomsen_univKK} and hence we have to verify two properties: That the composition is natural in $D$ and that it maps $1_{\Cz(X)}\in\KK_0^\Gamma(\Cz(X),\Cz(X))$ to $1_{\Cz(X)}\in\EE_0^\Gamma(\Cz(X),\Cz(X))$. Naturality is clear from the description of the three isomorphisms above and the definition of~$\Delta$.

For the verification of the other property, recall that $1_{\Cz(X)}\in\KK_0^\Gamma(\Cz(X),\Cz(X))$ is represented by the cycle $(\Cz(X)\oplus 0,0,\varepsilon'\oplus 0,\pi'\oplus 0)$, where $\varepsilon'$ and $\pi'$ denote the canonical representations of $\Gamma$ and $\Cz(X)$. After adding the degenerate cycle $\left(\HilbertMod_{\Cz(X)}\oplus\HilbertMod_{\Cz(X)}, \left(\begin{smallmatrix}0&\id\\\id&0\end{smallmatrix}\right)\!,\epsilon\oplus\epsilon,\pi\oplus\pi\right)$, we see that it is also represented by the cycle $\left(\HilbertMod_{\Cz(X)}\oplus\HilbertMod_{\Cz(X)}, \left(\begin{smallmatrix}0&U^*\\U&0\end{smallmatrix}\right)\!,\epsilon\oplus\epsilon,\pi\oplus\pi\right)$ for a~projection $U\colon \HilbertMod_{\Cz(X)}\cong\Cz(X)\oplus\HilbertMod_{\Cz(X)}\to\HilbertMod_{\Cz(X)}$ onto the second summand. Since $U$ is a module homomorphism, it has propagation zero and $U\in \PseudoLocLoc(X;\Cz(X))$. Furthermore we have $UU^*=\id$ and $p\coloneqq \id-U^*U$ is a projection onto an isomorphically embedded copy of $\Cz(X)$ in $\HilbertMod_{\Cz(X)}$.
For the calculation of the image of this class in $\EE_0^\Gamma(\Cz(X),D)$ we can choose the constant function $t\mapsto U_t\coloneqq U$. We end up with the element that is represented by the $*$-homomorphism $\Cz(X)\to \Kom(\HilbertMod_{\Cz(X)})$, $f\mapsto \pi(f)p$, and this is exactly $1_{\Cz(X)}$.

The isomorphism in odd degrees now follows from the one in even degrees by replacing $D$ with the suspension $D\otimes \Cz(\R)$.
\end{proof}
\begin{thebibliography}{99}
\bibitem[14]{EngelWulffZeidler}
Engel A., Wulff C., Zeidler R., Slant products on the {H}igson--{R}oe exact
 sequence, \href{https://doi.org/10.5802/aif.3406}{\textit{Ann. Inst. Fourier (Grenoble)}} \textbf{71} (2021),
 913--1021, \href{https://arxiv.org/abs/1909.03777}{arXiv:1909.03777}.

\bibitem[15]{GueHigTro}
Guentner E., Higson N., Trout J., Equivariant {$E$}-theory for
 {$C^*$}-algebras, \href{https://doi.org/10.1090/memo/0703}{\textit{Mem. Amer. Math. Soc.}} \textbf{148} (2000),
 viii+86~pages.

\bibitem[16]{GueWillYu_DyncomplKtheory}
Guentner E., Willett R., Yu G., Dynamical complexity and controlled operator
 {$K$}-theory, \href{https://arxiv.org/abs/1609.02093}{arXiv:1609.02093}.

\bibitem[17]{HankePapeSchick}
Hanke B., Pape D., Schick T., Codimension two index obstructions to positive
 scalar curvature, \href{https://doi.org/10.5802/aif.3000}{\textit{Ann. Inst. Fourier (Grenoble)}} \textbf{65} (2015),
 2681--2710, \href{https://arxiv.org/abs/1402.4094}{arXiv:1402.4094}.

\bibitem[19]{HigsonPedersenRoe}
Higson N., Pedersen E.K., Roe J., {$C^\ast$}-algebras and controlled topology,
 \href{https://doi.org/10.1023/A:1007705726771}{\textit{$K$-Theory}} \textbf{11} (1997), 209--239.

\bibitem[22]{HigRoe}
Higson N., Roe J., Analytic {$K$}-homology, \textit{Oxford Mathematical Monographs},
 Oxford University Press, Oxford, 2000.

\bibitem[32]{QiaoRoe}
Qiao Y., Roe J., On the localization algebra of {G}uoliang {Y}u, \href{https://doi.org/10.1515/FORUM.2010.036}{\textit{Forum
 Math.}} \textbf{22} (2010), 657--665.

\bibitem[37]{RoeCoarseGeometry}
Roe J., Lectures on coarse geometry, \textit{University Lecture Series},
 Vol.~31, \href{https://doi.org/10.1090/ulect/031}{Amer. Math. Soc.}, Providence, RI, 2003.

\bibitem[39]{Thomsen_univKK}
Thomsen K., The universal property of equivariant {$KK$}-theory,
 \href{https://doi.org/10.1515/crll.1998.112}{\textit{J.~Reine Angew. Math.}} \textbf{504} (1998), 55--71.

\bibitem[41]{WillettYuHigherIndexTheory}
Willett R., Yu G., Higher index theory, \textit{Cambridge Studies in Advanced
 Mathematics}, Vol.~189, \href{https://doi.org/10.1017/9781108867351}{Cambridge University Press}, Cambridge, 2020.

\bibitem[44]{WulffTwisted}
Wulff C., Coarse indices of twisted operators, \href{https://doi.org/10.1142/s179352531950033x}{\textit{J.~Topol. Anal.}}
 \textbf{11} (2019), 823--873, \href{https://arxiv.org/abs/1606.01297}{arXiv:1606.01297}.

\bibitem[47]{YuLocalization}
Yu G., Localization algebras and the coarse {B}aum--{C}onnes conjecture,
 \href{https://doi.org/10.1023/a:1007766031161}{\textit{$K$-Theory}} \textbf{11} (1997), 307--318.

\end{thebibliography}
\end{document}
